% ============================================================
%  深圳大学实验报告模板（SZU Experiment Report Template）
% This template can be modified according to the specific requirements of your experiment or project.
% Produced by: Chao Fan and Kewei Ou, AI School of SZU
% ============================================================

\documentclass[a4paper,12pt]{article}

% ----------------------------
% 中文与字体支持
% ----------------------------
\usepackage{fontspec}    % 字体支持
\usepackage{xeCJK}       % 中文字体支持
% 优先使用模板原字体；本机未安装时使用 TeX 发行版自带的 Fandol 宋体。
\IfFontExistsTF{Noto Serif CJK SC}{%
    \setCJKmainfont{Noto Serif CJK SC}%
}{%
    \setCJKmainfont{FandolSong-Regular.otf}[BoldFont=FandolSong-Bold.otf]%
}

% ----------------------------
% 页面设置与常用宏包
% ----------------------------
\usepackage{geometry}    % 页面边距控制
\geometry{left=1in, right=1in, top=1in, bottom=1in}

\usepackage{amsmath}
\usepackage{booktabs}
\usepackage{longtable}   % 支持长表格
\usepackage{graphicx}    % 插入图片
\usepackage{fancyhdr}    % 页眉页脚控制
\usepackage{tikz}        % 绘制框线等图形
\usetikzlibrary{calc}    % 坐标计算
\usepackage{verbatim}    % 显示代码块
\usepackage{float}       % 控制图片浮动位置（[H]参数）

% ----------------------------
% 页眉页脚设置
% ----------------------------
% \pagestyle{fancy}
\fancyhf{} % 清空默认页眉页脚
% \fancyhead[L]{深圳大学实验报告} % 左侧页眉文字
\fancyhead[C]{} % 中间空
\fancyhead[R]{} % 右侧空

% ============================================================
%                      文档开始
% ============================================================

\begin{document}

% ============================================================
% 封面页
% ============================================================
\begin{titlepage}
    \centering
    \vspace*{2cm}
    {\Huge\textbf{深 \ 圳 \ 大 \ 学 \ 实 \ 验 \ 报 \ 告}\par}
    \vspace{1.5cm}

    % 统一标签宽度、列间距与填写区宽度，使所有行的左右边界对齐。
    % \strut 保证不同内容的下划线与文字基线一致。
    \newcommand{\coverlabel}[1]{\makebox[3.2cm][s]{\textbf{#1}}}
    \newcommand{\covervalue}[1]{\underline{\makebox[8cm][c]{\strut #1}}}
    {\Large
    \begin{tabular}{@{}l@{\hspace{0.4cm}}l@{}}
        \coverlabel{课\hfill 程\hfill 名\hfill 称} &
        \covervalue{Python程序设计} \\[0.5cm]
        \coverlabel{项\hfill 目\hfill 名\hfill 称} &
        \covervalue{小项目一} \\[0.5cm]
        \coverlabel{学\hfill 院} &
        \covervalue{人工智能学院} \\[0.5cm]
        \coverlabel{专\hfill 业} &
        \covervalue{人工智能（IEEE荣誉班）} \\[0.5cm]
        \coverlabel{指\hfill 导\hfill 教\hfill 师} &
        \covervalue{舒婷} \\[0.5cm]
        \coverlabel{报\hfill 告\hfill 人} &
        \makebox[8cm][l]{%
            \underline{\makebox[3cm][c]{\strut 胡楷}}%
            \hspace{0.4cm}%
            \makebox[1.2cm][c]{\textbf{学号}}%
            \underline{\makebox[3.4cm][c]{\strut 2025681097}}%
        } \\[0.5cm]
        \coverlabel{实\hfill 验\hfill 时\hfill 间} &
        \covervalue{2026/11/18} \\[0.5cm]
        \coverlabel{提\hfill 交\hfill 时\hfill 间} &
        \covervalue{2026/11/22}
    \end{tabular}\par
    }

    \vfill
    {\Large\textbf{教务处制}\par}
\end{titlepage}

\newpage

% ============================================================

% 正文依据项目实现、运行日志和实际学习讨论扩写。
% 封面信息按此前要求保留；验证记录使用实际运行日期。
\setlength{\parindent}{2em}
\setlength{\parskip}{3pt}

\section{任务理解与实验准备}
\subsection{实验目标与完成标准}
这次作业是在给定框架中补齐 OpenCLIP ViT-B-32 图像编码器的六组 TODO，
使用 NumPy 完成从图片到特征向量的前向计算。模型参数由课程提供并保持固定，
因此实验重点是正确组织数组和实现计算过程，不涉及重新训练模型。
输入是预处理后的 float32 数组 $[B,3,224,224]$，输出是 $[B,512]$ 的单位向量。
这里 $B$ 表示批量图片数；内部的 $N=50$ 表示 token 数，$D=768$ 表示特征维度。

我对任务的理解是：程序需要根据输入图片实际计算，既要符合输入输出约定，
也要与教师参考结果在指定误差内一致。仅能生成一个形状正确的数组并不足以说明实现正确。
最终向量可以用于后续相似度比较，但本实验没有加入分类器，也不以分类准确率作为验收指标。

\noindent\textbf{代码仓库：}\texttt{https://github.com/mmga-137/project1-numpy}

\subsection{计算流程与实现范围}
图像先切成 49 个小块，经投影后加入 CLS 和位置编码，形成 50 个 token。
经过前置 LayerNorm 和 12 个 Transformer 块后，选择 CLS、执行后置 LayerNorm，
再投影到 512 维并进行单位化。六组补全内容及其检查位置如下。
\begin{center}
\small
\begin{tabular}{@{}lll@{}}
\toprule
任务 & 主要工作 & 对应验证\\
\midrule
TODO-01 & 图像切块与元素顺序保持 & patches\\
TODO-02 & 将特征维拆成多个注意力头 & 单元测试、block0\\
TODO-03 & 注意力分数、加权汇总与合并 & block0\\
TODO-04 & MLP 与两次残差相加 & block0、last\_tokens\\
TODO-05 & 块投影、CLS 拼接和位置编码 & tokens\\
TODO-06 & 选择 CLS、投影与最终输出 & output、batch\\
\bottomrule
\end{tabular}
\end{center}
图像预处理、LayerNorm、精确 GELU、参数加载及教师检查脚本沿用课程文件。
模型内部宽度为 768，采用 12 个头，每头 64 维；MLP 中间宽度为 3072。

\subsection{环境与复现步骤}
已验证环境为 Windows 11、Python 3.13.12、NumPy 2.4.2、Pillow 12.1.1。
课程附件包含 152 个视觉参数张量、十张本地图片和阶段参考数组。
从原附件恢复完整的 \texttt{weights/} 与 \texttt{assets/} 后，在仓库根目录执行：
\begin{verbatim}
python -m pip install -r requirements.txt
python check.py --stage all
python run.py --output output_tokens.npz
python -m unittest discover -s tests -v
\end{verbatim}
数据与参数的恢复办法写在 README 中，大型课程资源不重复上传。
基础验证日期为 2026/11/24（UTC+8），见 \texttt{results/environment.json}；
封面日期保留此前指定值，文档修订与测试时间分别记录。

\clearpage
\section{从像素到图片块特征}
\subsection{预处理与输入检查}
课程预处理先将图片最短边双三次缩放至 224，再中心裁剪为 $224\times224$，
转换成 RGB。像素转为 float32 并除以 255，随后调整为 CHW 排列，
按给定通道均值和标准差归一化，最后增加 batch 轴。
同样的颜色值如果预处理顺序或归一化参数不同，后续特征也可能不同，因此沿用给定实现。

正式切块前还检查输入是否为 NumPy 数组、形状是否严格为 $[B,3,224,224]$、
类型是否为 float32。形状不符抛出 ValueError，类型不符抛出 TypeError。
该检查能尽早提示尺寸或通道顺序错误，避免错误数据进入后续计算。

\subsection{切块时为什么需要轴重排}
每个块的边长为 32，故图片每边有 $224/32=7$ 块，共 $7\times7=49$ 块。
每块包括三个通道，按 C、H、W 展开后是 $3\times32\times32=3072$ 个数。
切块前后的形状为 $[B,3,224,224]\to[B,49,3072]$，具体分三步：
\[
\begin{aligned}
[B,3,224,224]&\longrightarrow[B,3,7,32,7,32]\\
&\longrightarrow[B,7,7,3,32,32]\\
&\longrightarrow[B,49,3072].
\end{aligned}
\]
第一次 reshape 将高、宽分别拆为网格轴与块内轴；
\texttt{transpose(0,2,4,1,3,5)} 将同一块的数据放到一起；最后合并网格轴和块内轴。
网格按从左到右、再从上到下排列，块内按通道、行、列展开。

我用 $4\times4$ 的单通道矩阵理解这一点。若像素从 0 到 15 逐行排列，
左上角 $2\times2$ 块应展开为 $[0,1,4,5]$，右上角为 $[2,3,6,7]$。
直接把连续数据每四个分组，第一组却是 $[0,1,2,3]$，对应原图第一行。
因此，reshape 能改变形状，却不会自动理解哪个方形区域是一块。

\subsection{块投影、CLS 与位置编码}
块投影使用给定权重，对每块的全部 3072 个数进行加权组合，输出 768 个特征：
\[
E=P W_p^{\mathsf T},\qquad
[B,49,3072]\ @\ [3072,768]\longrightarrow[B,49,768].
\]
这里没有偏置，49 个块使用同一套投影参数。
每个输出特征由整个块的输入加权得到，不能把它理解为每四个像素合成一个值。
权重存为 $[768,3072]$，因此计算时需要转置。

先将 CLS 复制到每张图片，并在 49 个块向量前拼接，得到 50 个 token。
此时形状为 $[B,50,768]$，其中 CLS 用于在后续注意力中汇总全图信息。
位置参数为 $[50,768]$，通过广播加到每张图片上，让模型能够区分各 token 的位置。
进入编码块之前加入位置编码，后续特征同时利用内容与位置信息。

\clearpage
\section{Transformer 与最终输出}
\subsection{LayerNorm 与多头注意力}
每个 Transformer 块采用前置 LayerNorm。对单个 token 的特征计算均值和方差，
再标准化，并施加模型给定的缩放 $\gamma$ 和偏置 $\beta$：
\[
\operatorname{LN}(x)=\frac{x-\mu}{\sqrt{\sigma^2+10^{-5}}}\odot\gamma+\beta.
\]
它调整每个 token 内部的数值尺度，保持数组形状不变；最终输出不一定严格具有零均值和单位方差。
标准化有助于控制后续计算尺度，稳定项 $10^{-5}$ 避免方差为零时除零。

注意力先以线性变换得到 Q、K、V，分别用于发起匹配、参与匹配和提供被汇总的信息。
三者均从 $[B,N,768]$ 拆为 $[B,N,12,64]$，再转置为 $[B,12,N,64]$。
每个头使用一组 64 维特征，但仍能关注全部 50 个 token。
对每个头计算：
\[
S=\frac{QK^{\mathsf T}}{\sqrt{64}},\qquad
A=\operatorname{softmax}(S),\qquad C=AV.
\]
分数和系数的形状为 $[B,12,N,N]$。缩放因子控制点积分数的尺度；
softmax 沿最后一轴转换为非负、和为 1 的权重。
给定辅助函数先减去每行最大值再取指数，改善数值稳定性。
实际 K 为四维，使用 \texttt{swapaxes(-1,-2)} 只交换末尾两个轴。

加权后的结果从 $[B,12,N,64]$ 转回 $[B,N,12,64]$，合并成 $[B,N,768]$，
再做输出投影。图像使用 \texttt{causal=False}，允许各位置互相关注；
接口保留 causal 开关，并用合成参数检查了该分支。

\subsection{两次残差连接与 MLP}
注意力提供 token 之间的信息交互，MLP 对每个 token 内部的特征进行非线性加工。
两次更新分别加回对应的输入，完整计算为：
\[
\begin{aligned}
R&=X+\operatorname{Attention}(\operatorname{LN}_1(X)),\\
U&=\operatorname{LN}_2(R),\\
M&=\operatorname{GELU}(UW_1^{\mathsf T}+b_1)W_2^{\mathsf T}+b_2,\\
Y&=R+M.
\end{aligned}
\]
第一条残差加回原始 $X$，第二条加回 $R$，为原特征提供直接传递路径。
MLP 将每个 token 从 768 维扩展至 3072 维，经 GELU 后再投影回 768 维，token 数量始终不变。
GELU 使用给定的精确 erf 版本，逐元素进行平滑的非线性变换；
如果没有非线性激活，连续两层线性变换仍可合并为一次。
普通线性权重按 $[\mathrm{out},\mathrm{in}]$ 存储，故使用 \texttt{weight.T}。

\subsection{全图特征的取得}
12 个块结构相同、参数分别独立，内部序列最终仍为 $[B,50,768]$。
通过 \texttt{tokens[:,0,:]} 选择 CLS，后置 LayerNorm 后乘 $[768,512]$ 的
\texttt{visual.proj}。该矩阵已经按输入、输出维排列，不再转置。
最后调用单位化函数，将每行除以其 L2 长度，得到 $[B,512]$。
这一步控制整条输出向量的长度，与前面 LayerNorm 调整特征分布的作用不同。

\clearpage
\section{结果验证与问题分析}
\subsection{逐阶段与教师参考结果比较}
我判断是否完成时，还需要看数值是否符合要求。这里将本程序结果 $a$ 与教师参考 $r$ 逐元素比较：
\[
 e=|a-r|,\qquad t=\mathrm{atol}+\mathrm{rtol}|r|,\qquad \rho=e/t.
\]
其中 $e$ 是绝对误差，$t$ 是允许误差。所有阶段的 $\mathrm{rtol}=2\times10^{-4}$；
表中最大 $\rho\leq1$，才说明每个元素都符合范围。
\begin{center}
\small
\setlength{\tabcolsep}{4pt}
\begin{tabular}{@{}lrrrrc@{}}
\toprule
阶段 & 最大绝对误差 & 平均绝对误差 & atol & 最大 $\rho$ & 判定\\
\midrule
preprocess & 0 & 0 & $2.000\times10^{-6}$ & 0.0000 & 通过\\
patches & 0 & 0 & $2.000\times10^{-4}$ & 0.0000 & 通过\\
tokens & $4.387\times10^{-5}$ & $1.335\times10^{-6}$ & $2.000\times10^{-4}$ & 0.0769 & 通过\\
block0 & $7.147\times10^{-6}$ & $3.348\times10^{-7}$ & $2.000\times10^{-4}$ & 0.0351 & 通过\\
last\_tokens & $4.196\times10^{-5}$ & $2.633\times10^{-6}$ & $2.000\times10^{-4}$ & 0.0948 & 通过\\
output & $5.662\times10^{-7}$ & $8.688\times10^{-8}$ & $2.000\times10^{-5}$ & 0.0205 & 通过\\
batch & $4.061\times10^{-7}$ & $7.848\times10^{-8}$ & $2.000\times10^{-5}$ & 0.0145 & 通过\\
\bottomrule
\end{tabular}
\end{center}
assets 的完整性检查也通过，因此八个阶段全部通过。
最终输出的 5120 个元素均未超限，最大 $\rho=0.02053$。

\subsection{具体特征值的对照}
下面选取首值、绝对误差最大值、参考值最接近零的位置及末值作直接对照。
图号对应课程图片，下标从 0 开始；展示保留 9 位小数，误差由未舍入数据计算。
\begin{center}
\small
\setlength{\tabcolsep}{3pt}
\begin{tabular}{@{}lrrrr@{}}
\toprule
图号/下标 & 本程序输出 & 教师参考 & 绝对误差 & 允许误差\\
\midrule
01/0 & -0.039898202 & -0.039898291 & $8.941\times10^{-8}$ & $2.798\times10^{-5}$\\
10/152 & 0.113599643 & 0.113600209 & $5.662\times10^{-7}$ & $4.272\times10^{-5}$\\
05/0 & -0.000007132 & -0.000007156 & $2.465\times10^{-8}$ & $2.000\times10^{-5}$\\
10/511 & -0.000589774 & -0.000589705 & $6.880\times10^{-8}$ & $2.012\times10^{-5}$\\
\bottomrule
\end{tabular}
\end{center}
图 10 下标 152 处的最大误差为 $5.662\times10^{-7}$，小于允许的 $4.272\times10^{-5}$。
我理解，数值不必完全相同，但必须逐个符合教师给定的范围。
输出为有限 float32 数组 $(10,512)$，L2 长度介于 $0.99999988$ 与 $1.00000000$ 之间。

\subsection{补充验证与问题修复}
六个单元测试及换序、重复图片的 batch 检查通过。
新输入单张与四张批量输出最大差为 $1.490\times10^{-7}$；
没有 \texttt{assets/} 的独立项目完成 RGB、灰度新图推理，命令行与 API 输出最大差为 0，
说明推理不需要教师参考数组。

原切块会接受 $(1,3,112,448)$ 的错误尺寸，现已检查完整形状与 float32 类型。
切块顺序诊断中，直接 reshape 有 150464 个元素错位，正式实现为 0。
完整数值对照见 \texttt{results/reference\_comparison.json}，复现方法见 README。

\clearpage

\section{学习反思与后续改进}
\subsection{从能看懂形状，到能解释计算}
对我这个大二学生来说，这次作业最有价值的地方，是把“图像特征提取”拆成了可以逐步解释的数组运算。
刚开始看到 $3072\to768$，我以为维度减为四分之一，就是把数据“除以 4”，
后来又把它想成每个 $2\times2$ 小区域合并出一个数。
经过小例子推演，我才分清：3072 是一个图片块的全部像素数，768 是模型规定的输出特征数；
每个输出特征都由整个块的输入加权计算，768 组权重产生 768 个结果。
这个误解让我意识到，知道两个维度相差多少，还不等于知道程序在算什么。

切块也是类似的情况。$[B,49,3072]$ 看起来很整齐，但形状正确不能保证像素放在正确位置。
我能从小矩阵中找出右上角的 $[2,3,6,7]$，再看 transpose 的轴顺序时，
就能把它理解为“把同一块的数据安排到一起”。以后遇到多维数组，我会先写清每个轴的含义，
用一个小例子检查顺序，再考虑正式模型的尺寸。

\subsection{对 Transformer 的理解更具体了}
我一开始把 Transformer 的作用理解成“得到带位置的信息”，后来才明确位置编码是在进入它之前就加上的。
块投影负责得到局部特征，位置编码补充位置信息，注意力再根据相关程度交换信息，
CLS 经过多层更新后用于汇总整张图片。这些步骤承担的任务不同，不能只记一个模块名称。

对于“12 个头”，我需要记清它拆分的是特征维，每个头仍能关注全部 50 个 token。
MLP 中的 $768\to3072\to768$ 同样改变特征维度，不改变 token 数量。
两次残差相加的小例子则帮助我区分第二次要加回第一次更新后的结果。
GELU 提供非线性加工，LayerNorm 调整数值尺度，它们不负责增加图片块的数量。

\subsection{如何判断完成，以及还需要补什么}
这次作业让我开始区分“运行成功”“数值正确”和“自己理解了”这三件事。
八个阶段通过、有新输入和独立运行记录，可以支持当前实现满足给定模型的要求；
但若我不能解释轴顺序、权重方向或错误信息，就还不能只凭 PASS 认为自己已经掌握。
我希望之后把“每一步输入是什么、输出是什么、为什么这样算”作为检查理解的方法。

目前的参考图片数量有限，额外合成输入主要验证程序行为，并没有测量分类或检索效果。
后续可以用更多自选图片开展相似度检索，明确评价指标后再判断实际效果。
性能方面，给定 GELU 的逐元素 erf 循环值得进一步研究；若采用等价向量化实现，
需要重新验证精度并测量耗时，不能直接把设想写成已经取得的提升。
我也还需要通过继续阅读和练习，理解这些权重在训练时是怎样得到的。

\section{本人参与与辅助说明}
这次作业中，我确实使用了 AI 辅助学习。我先提出不理解的地方，与 AI 讨论，
再自己推演切块和残差的小例子，把我的理解反馈给 AI 核对。
本人还参与任务资料整理、报告要求确定、GitHub 仓库创建与复查。
代码实现、数值检查、报告整理及最终 GitHub 提交也有 AI 协助。
辅助工具：OpenAI Codex。教师提供原始框架、参数、图片、参考结果和数值辅助代码。
结论以仓库中的可复现记录为依据，反思按本次实际讨论与理解过程整理。

\clearpage

% ============================================================
% 批阅与成绩评定页

% ============================================================

% 绘制正文外框（包含批阅区域）
\begin{tikzpicture}[remember picture, overlay]
  \draw[thick]
    ($(current page.north west)+(2cm,-3cm)$)
    rectangle
    ($(current page.south east)+(-2cm,3.5cm)$);
\end{tikzpicture}

\vspace{1cm}

% 批阅区
\noindent \textbf{指导教师批阅意见：}
\vspace{6cm}
\hfill

\vspace{1cm}

\noindent \textbf{成绩评定：}
\vspace{2cm}
\hfill

\vspace{1cm}

\noindent
\hfill
\textbf{指导教师签字：}\quad\quad\quad\quad\quad\\[0.5cm]

\hfill
\textbf{年}\quad\quad
\textbf{月}\quad\quad
\textbf{日}\quad

\vspace{1cm}

% 备注部分
\noindent \textbf{备注：}

% 框外的注释
\begin{tikzpicture}[remember picture, overlay]
  \node[
    anchor=north west,
    inner sep=0pt,
    text width=\dimexpr\paperwidth-4cm\relax
  ] at ($(current page.south west)+(2cm,3.1cm)$) {
    \begin{minipage}{\dimexpr\paperwidth-4cm\relax}
      \small
      \noindent
      \textbf{注：}\;1、报告内的项目或内容设置，可根据实际情况加以调整和补充。\\[0.2cm]
      \phantom{\textbf{注：}\;}2、教师批改学生实验报告时间应在学生提交实验报告时间后 10 日内。
    \end{minipage}
  };
\end{tikzpicture}

% ============================================================
%                      文档结束
% ============================================================

\end{document}
